\section{Mechanism: Saturated Basin, Load-Bearing Residual}
\label{sec:convergence}

\S\ref{sec:saturation} reported that V-axis supervision moves the
class-mean component of the network's features by up to $+0.36$ while
test accuracy falls and training accuracy stays at ceiling. We take
the class-mean/residual decomposition as the operative geometry and
ask what each subspace actually does for the trained network.

\paragraph{The class-mean basin is saturated across seeds.}
Across the ten same-architecture checkpoints that make up our SOTA
ensemble (\S\ref{sec:sota}; EMOD $d{=}6$, five seeds $\times$ two
training lengths), the class-mean subspace encodes the CLIP-text V-axis at a
consistent strength ($|r|\in[0.60, 0.77]$, unsigned): every
well-trained seed of this architecture encodes the same class-mean
valence direction, all $|r|$ in a tight band and the same sign (the
nine-story Qwen3 axis more weakly, $|r|\sim0.3$--$0.5$, per-checkpoint
n.s.; Appendix~\ref{app:convergence-deep-dive}). This is the \emph{saturated
basin}---the component \S\ref{sec:saturation} showed auxiliary
supervision cannot improve. What differs across seeds, and what
carries the recoverable gain, is the orthogonal within-class residual.
(We do \emph{not} claim architecture-independent convergence: pooling
V-axis encoding against accuracy across heterogeneous backbones yields
an inflated correlation that collapses within any single architecture
and does not survive a matched-norm null; see
Appendix~\ref{app:convergence-deep-dive}.)

\begin{revblock}
\paragraph{Residual encoding and ensemble contribution.}
On $20$ independently trained SOTA-recipe checkpoints (one per seed),
residual encoding correlates with leave-one-out ensemble contribution when the
residual direction is estimated on test subjects ($r{=}0.51$, $p{=}0.022$) but not
when fit on training subjects only ($r{=}0.21$, $p{=}0.37$;
Appendix~\ref{app:convergence-deep-dive}). The link is suggestive; the causal
ablation below is the load-bearing evidence.

\paragraph{Direct test: ablate the residual at inference time.}
A directional ablation~\citep{arditi2024refusal}---the projection
$(I - vv^\top)$ with $v$ the unit CLIP-derived within-class V-axis
residual, fit on the $80$ training subjects only---drops test accuracy on
all $20$ independent checkpoints (mean $\Delta\mathrm{BACC}{=}{-}0.018$;
every drop lies below all $50$ matched random directions, mean
$z{\approx}{-}8.4$). The residual is genuinely load-bearing: ablating
$v$ orthogonalised to the class-mean subspace still lowers accuracy on
$18/20$ checkpoints (mean ${-}0.007$, $z{\approx}{-}4.0$). Along $v$, clip identity
explains $13.8\%$ of within-class test variance (random directions: $6.3\%$) and
subject identity $9.0\%$ ($14.6\%$; both $20/20$ checkpoints): the residual
mainly encodes \emph{which clip} is being watched, not the subject.
\end{revblock}
This
decomposition---saturated
basin, load-bearing residual---is the operative geometry that
\S\ref{sec:saturation} predicted from the failure of V-axis
supervision and that \S\ref{sec:sota} turns into a positive
prescription.
